
Differentiability Explorer

Project Overview

The Differentiability Explorer is an interactive web application developed using HTML, CSS, and JavaScript to investigate differentiability and continuity through numerical approximations.

The application allows users to select mathematical functions, adjust an evaluation point (a), and modify a step size (h) to observe how numerical derivative estimates behave.

The primary objective is to demonstrate how the left- and right-hand difference quotients can provide evidence of differentiability at a selected point.

Mathematical Background

A function is differentiable at a point (a) when the following limit exists and is finite:

[
f'(a)=\lim_{h\to0}\frac{f(a+h)-f(a)}{h}
]

The application approximates this derivative using two one-sided difference quotients.

Left-hand derivative estimate:

[
D_L(a,h)=\frac{f(a)-f(a-h)}{h}
]

Right-hand derivative estimate:

[
D_R(a,h)=\frac{f(a+h)-f(a)}{h}
]

As (h) approaches zero, these estimates should approach the same finite value if the function is differentiable at (a).

Differentiability also requires continuity. However, a continuous function may not be differentiable at every point, as demonstrated by the absolute value function at (a=0).

Application Features

Interactive function selection, including smooth and nondifferentiable examples.

Adjustable evaluation point (a) and step size (h).

Graphical representation of the selected function.

Numerical left- and right-hand derivative estimates.

Secant and tangent-line approximations.

Diagnostic information displaying function values and estimated slopes.

Numerical classification of likely differentiability, possible vertical tangents, or inconclusive results.

Theory and explanatory content to support mathematical interpretation.

Numerical Classification

The application compares the left- and right-hand derivative estimates using an absolute difference:

[
\Delta=|D_L-D_R|
]

A numerical tolerance is used to determine whether the estimates are sufficiently close to suggest differentiability.

The application also evaluates whether slopes are unusually large as a possible indication of a vertical tangent.

These classifications are numerical approximations rather than formal mathematical proofs. A small difference between the estimates does not guarantee differentiability, and a large slope does not necessarily imply a vertical tangent.

How to Run

Download or clone the repository.

Open the application's HTML file in a modern web browser.

Select a mathematical function.

Adjust the evaluation point (a) and step size (h).

Observe the graph, numerical diagnostics, and differentiability conclusion.

No additional software installation is required for the standalone HTML application.

Suggested Experiments

Experiment 1: Smooth function

Select (f(x)=x^2) and evaluate the function at (a=2). Decrease (h) and observe how the left- and right-hand estimates approach the exact derivative (f'(2)=4).

Experiment 2: Sharp corner

Select (f(x)=|x|) and set (a=0). Observe how the one-sided estimates disagree, indicating nondifferentiability.

Experiment 3: Vertical tangent

Select the application's vertical-tangent example and examine how the slope estimates change as (h) decreases.

Experiment 4: Piecewise function

Select a piecewise function and investigate its behavior near the boundary between its pieces.

Limitations

The application uses finite step sizes rather than evaluating exact limits. Consequently, numerical results are subject to approximation error.

Very small step sizes can introduce floating-point rounding errors and cancellation, while larger step sizes may produce less accurate derivative estimates.

The numerical tolerance and slope thresholds are heuristic values and may produce incorrect classifications for certain functions.

Therefore, results should be interpreted as numerical evidence rather than definitive mathematical conclusions.

Technologies Used

HTML: Application structure and interface

CSS: Visual layout and styling

JavaScript: Function evaluation, numerical calculations, interactive controls, and graph updates

Educational Purpose

This project was developed to explore the relationship between calculus, numerical analysis, and computational mathematics. It demonstrates how mathematical definitions can be translated into algorithms and how numerical approximations can help visualize abstract concepts.

The project also highlights the importance of numerical accuracy, mathematical assumptions, and critical interpretation of computational results.

AI Assistance Disclosure

ChatGPT was used as a learning and development resource during the project, including assistance with JavaScript implementation, debugging, mathematical explanations, numerical-method discussions, and documentation.

The application was developed and reviewed through an iterative process involving mathematical reasoning, testing, and revision. AI-generated suggestions were evaluated for mathematical accuracy and suitability for the project's objectives.
